% =====================================================================
% CHAPTER 4 — THE BASELINE STRATEGY (L1)  (~12% of total length)
%
% Job: answer RQ1 cleanly and build the reference point Chapters 5–6
% are measured against. Replication + honest evaluation — NOT the
% novelty chapter, and it shouldn't strain to be.
%
% Structure:
%   4.1 The threshold strategy (rule, tuning of L on train split)
%   4.2 Results vs baselines (buy-and-hold, price-only momentum)
%   4.3 Cost sensitivity
%   4.4 Answer to RQ1
%
% REMEMBER the success criterion from the plan: a working,
% leakage-controlled backtest with honest uncertainty. Profitability
% is NOT required. If the strategy doesn't beat buy-and-hold, say so
% plainly — L2/L3 do not depend on L1 being profitable.
% =====================================================================
\chapter{The Baseline Strategy (L1)}
\label{ch:layer1}

\section{The Threshold Strategy}
\label{sec:l1-strategy}

[The rule from \textcite{mitic2026llm}: trade when $|S_t| > L$, with
$L$ tuned on the training split only. Spell out the position
convention, holding period, and exact tuning grid.]

\section{Results}
\label{sec:l1-results}

% Every P/L table: with AND without costs, per the rigor protocol.
\begin{table}[htbp]
  \centering
  \caption{L1 backtest vs.\ baselines, out-of-sample. [Placeholder
  values. Report each index separately; add a costs panel.]}
  \label{tab:l1-results}
  \begin{tabular}{@{}l S[table-format=2.1] S[table-format=1.2] S[table-format=2.1] S[table-format=2.1]@{}}
    \toprule
    Strategy & {Cum.\ P/L (\%)} & {Sharpe} & {Hit rate (\%)} & {Max DD (\%)} \\
    \midrule
    Buy-and-hold        & 4.2 & 0.31 & {---} & 12.4 \\
    Price-only momentum & 3.1 & 0.24 & 51.2  & 10.8 \\
    Threshold on $S_t$  & 6.8 & 0.52 & 53.0  &  9.6 \\
    \bottomrule
  \end{tabular}
\end{table}

[State the comparison in words; give uncertainty (e.g.\ bootstrap CIs
on Sharpe differences) rather than a bare point estimate.]

\section{Cost Sensitivity}
\label{sec:l1-costs}

[P/L as a function of assumed round-trip cost; the break-even cost is
a more honest summary than any single net figure.]

\section{Answer to RQ1}
\label{sec:l1-answer}

% One short section closing the question — mirrors the RQ wording
% from Section 1.2. Repeat this pattern in Chapters 5 and 6.
[Direct answer, with the caveats that survive Chapter 3's protocol.]
